\section{1901 | The first Portuguese plates}
\subsection*{Motor car circulation licences}

The first plates appeared on 3 October 1901, then called “motor car circulation
licences”. They were issued by each district's Civil Government, and only one
plate, on the back of the vehicle, was required, although photographs of the
time show one was sometimes fitted at the front too.

Each combination was the name or abbreviation of the district seat that issued
the licence followed by a number, so plates such as \plate{LISBOA 130} or
\plate{LXA 256} could be seen.

The decree of that date required every car to carry on its back a metal plate
with the licence number and the district seat where it was granted.

The first car registration in Portugal was in the name of Augusto Sabido
Teixeira de Aragão, from Santiago do Cacém.

\section{1911 | The X-000 series}
\subsection*{First standardisation of the Portuguese plate system}

As the number of cars grew, the system was standardised nationwide. The country
was split into north, centre and south, and each zone's initial came before the
digits. Plates were black with white characters, with sequences such as
\plate{N-000}, \plate{C-000} and \plate{S-000}, or \plate{N-00000},
\plate{C-00000} and \plate{S-00000}.

\platefig{fmt-1911}{A 1911-series plate: \plate{S-26384}.}

\begin{itemize}
  \item Hire cars had an \plate{A} added after the serial number.
        E.g. \plate{C-4569-A}.
  \item Provisional plates had the letters \plate{WW} added.
        E.g. \plate{S-1703-WW}.
  \item Army cars used the \plate{MG-000} sequence, which stayed in use until
        the 1950s.
\end{itemize}

\section{1918 | Registration zones for the Azores and Madeira}

The Azores and Madeira were given the \plate{A-000} and \plate{M-000} series,
respectively. Unlike on the mainland, this system remained in force on the
islands until 1962.

\section{1937 | The XX-00-00 series}
\subsection*{Second standardisation of the Portuguese plate system}

In some zones the serial number had already reached five digits, which made
vehicles hard to identify at a glance. This is when the familiar combinations of
two letters and two pairs of digits began.

On 1 January 1937 the second national registration system came into force. The
division into zones was kept, but now by letter groups:
\begin{itemize}
  \item \plate{AA} to \plate{LZ} for the Lisbon zone;
  \item \plate{MA} to \plate{TZ} for the Porto zone;
  \item \plate{UA} to \plate{ZZ} for the Coimbra zone (the former Centre zone).
\end{itemize}

Plates stayed black with white characters, except for diplomatic vehicles and
vehicles under temporary import, and were made of plastic with raised
characters.

\platefig{1937-pg0870}{A 1937-series plate from the Porto zone: \plate{PG-08-70}.}

The letters \plate{Q} and \plate{J} were not allowed because they could be
mistaken for \plate{O} and \plate{I}, and the letter groups \plate{CU} and
\plate{FD} were banned for reasons of “decency”.

Some letter groups were reserved for organisations and bodies:\footnote{According to Portuguese Wikipedia (\emph{Matrículas de automóveis em Portugal}); the 2023 edition did not mention them.}
\begin{itemize}
  \item \plate{MG}, \plate{ME} and \plate{MX} for the Army, \plate{AP} for the
        Navy, \plate{AM} for the Air Force and \plate{EP} for the State civil
        administration;
  \item \plate{CD}, \plate{CC} and \plate{FM}, in red on white, for the Diplomatic
        Corps, the Consular Corps and other officials on diplomatic missions
        (until 1985, see below);
  \item \plate{TA}, \plate{TB} and \plate{TC}, in white on red, for vehicles
        under temporary import in the Lisbon, Porto and Coimbra zones.
\end{itemize}
The National Republican Guard and the former Fiscal Guard had their own
sequences, starting with \plate{GNR} and \plate{GF} respectively.

\subsection{Converting the 1911 plates}

Cars registered under the 1911 system had to move to the new one, and were given
specific series according to their old plate. The last four digits of the old
number became the two pairs of digits.

\begin{itemize}
  \item \textbf{Lisbon:} \plate{S-000} and \plate{S-0000} became \plate{AA-00-00};
        \plate{S-10000} became \plate{AB-00-00}; \plate{S-20000},
        \plate{AC-00-00}; \plate{S-30000}, \plate{AD-00-00}.
        E.g. \plate{S-12345} became \plate{AB-23-45}.
  \item \textbf{Porto:} \plate{N-000} and \plate{N-0000} became \plate{MM-00-00};
        \plate{N-10000} became \plate{MN-00-00}.
        E.g. \plate{N-1234} became \plate{MM-12-34}.
  \item \textbf{Coimbra:} every vehicle moved to the \plate{UU-00-00} series.
        E.g. \plate{C-123} became \plate{UU-01-23}.
\end{itemize}

\section{1962 | The islands adopt the system}

The Azores and Madeira only adopted the new system in 1962. The Azores were split
into three registration zones, matching the Autonomous Districts of Ponta
Delgada, Angra do Heroísmo and Horta, each with its own registration office, and
Madeira was a single zone, matching the Autonomous District of Funchal. The
letters reserved for these zones were:
\begin{itemize}
  \item \plate{AN} for Angra do Heroísmo;
  \item \plate{AR} and \plate{AS} for Ponta Delgada;
  \item \plate{HO} for Horta;
  \item \plate{MA} and \plate{MD} for Funchal.
\end{itemize}

The letters \plate{TD}, \plate{TG}, \plate{TH} and \plate{TF}, in white on a red
background, were reserved for vehicles under temporary import in Ponta Delgada,
Angra do Heroísmo, Horta and Funchal, respectively.

\section{1970s | New regional offices}

During the 1970s new regional offices and registration zones were set up, each
with reserved letter groups:
\begin{itemize}
  \item \plate{EM} and \plate{EV} for Évora;
  \item \plate{ZA} and \plate{ZB} for Braga;
  \item \plate{ZC} and \plate{ZD} for Vila Real;
  \item \plate{ZE} and \plate{ZF} for Aveiro;
  \item \plate{ZG} and \plate{ZH} for Guarda;
  \item \plate{ZI} and \plate{ZL} for Santarém;
  \item \plate{ZM} and \plate{ZN} for Setúbal;
  \item \plate{ZO} and \plate{ZP} for Faro.
\end{itemize}

The Évora office was considered the regional hub, so it was also given a code for
vehicles under temporary import: \plate{TE}. Because registration by zone was
dropped shortly afterwards, most of these letter groups were never actually used
to identify a zone.

When the Portuguese colonies became independent, anyone who wished could import
their car to Portugal. New plates were issued for this, starting with
\plate{VZ}, \plate{ZO}, \plate{ZE} and \plate{UZ}.\footnote{It is unconfirmed whether \plate{ZE} and \plate{UZ} were used for these imports.}
\plate{VZ} was jokingly read as “Veio Zangado” (“came back angry”).

\section{1980s | Extending the 1937 system}

In the early 1980s, as the series reserved for some zones ran out, the zones were
abolished and vehicles were registered sequentially nationwide.

Even so, the \plate{AA-00-00} series was quickly nearing its end. To delay that,
the “famous” \plate{CU} and \plate{FD} combinations and the letters \plate{Q} and
\plate{J} were allowed.

In 1985 diplomatic plates moved to their own sequence, in the formats
\plate{000-CD000}, \plate{000-CC000} and \plate{000-FM000}. The first group of
digits identifies the country of the vehicle's user and the letters its
diplomatic status (see section~\ref{sec:diplomaticos}).

\platefig{diplomatic-fm}{Diplomatic plate \plate{024-FM256}.}

\section{1992 | The 00-00-XX series}
\subsection*{New numbering and the euroband}

In 1992 plates changed both look and numbering. White reflective plates with
black characters and a blue band on the left (the euroband), with the European
Union emblem and the country's initial, \plate{P}, came into use. The typeface
is DIN 1451, created by the German Institute for Standardisation in 1931. The new
plates were made of metal, and Portugal was the second country to adopt the
euroband.

The letters moved to the end: \plate{00-00-AA}. The two changes did not happen at
the same time: the new plate applies to registrations from 1 January 1992, but
\plate{00-00-AA} numbering only began on 1 March (Decree-Law 54/2005, arts.~3
and~5). So in the first two months of 1992, \plate{AA-00-00} numbers were issued
on the new white plates.

In the new sequence no letter group was reserved for special bodies. The only
exception was the 1997 imports (section~\ref{sec:importados}).

\platefig{1992-7293fq}{A 1992-series plate without the yellow band: \plate{72-93-FQ}.}

\section{1998 | Year and month of first registration}

The 1992 numbering was kept, but the look changed: a yellow band on the right now
showed the year and month of the vehicle's first registration.

\platefig{1998-4572xq}{\plate{45-72-XQ}, with the yellow band for June 2004.}

\section{2005 | The 00-XX-00 series}

With every combination of the previous numbering used up, plates moved to the
\plate{00-XX-00} format in 2005.

\platefig{2005-99sa99}{A 2005-series plate: \plate{99-SA-99}.}

\section{2020 | The XX 00 XX series}

In 2020 both the numbering and the look changed, to the system in use today.
Decree-Law 2/2020 of 14 January defined the new number as “two groups of two
letters and a central group of two digits”, a format Decree-Law 54/2005 had
already provided for once the \plate{00-AA-00} series ran out. According to IMT,
it applies only to new vehicles; the owner of a vehicle from an earlier series
can ask for a plate without the month and year. The first plate, \plate{AA 01 AA},
was issued on 3 March 2020 to an electric car; IMT reserved \plate{AA 00 AA}
“for posterity”.

\platefig{2020-aa16aa}{A 2020-series plate: \plate{AA 16 AA}.}

The dashes were replaced by spaces, because there would not be enough room for
numbers with several wide letters. The yellow band with the year and month of
first registration was removed to match the plates of the other European Union
countries: Portugal was the only one of the 27 member states with this
information, which could be mistaken for the plate validity date some countries
show.

\subsection{The vowel rule}

With four letters, many combinations are possible and some may be inappropriate.
So, on plates in this format issued by IMT's computer system, vowels are not used
in the second or sixth position. In IMT's words:

\begin{quote}\small
“não são utilizadas vogais na segunda e sexta posição desses caracteres nas
matrículas do novo formato […] Tal não se aplica, quando qualquer dos grupos de
duas letras incluir duas vogais iguais.”\\
(“vowels are not used in the second and sixth positions of new-format plates
[…] This does not apply when either two-letter group contains two identical
vowels.”)
\end{quote}

So \plate{AB 31 CE} cannot exist (a vowel in the 6th position), but
\plate{AA 12 BE} can, because the \plate{AA} group lifts the rule. IMT also says
it avoids “letter combinations that could form certain words”, without
publishing the list.

\subsection{The letters K, W and Y}

Until 2020 the letters \plate{K}, \plate{W} and \plate{Y} were not used on
ordinary plates (except \plate{K} for the 1997–1998 imports). Following the
Portuguese spelling reform, IMT announced they would become part of the new
series; we found no confirmation that any have been issued yet.
